\documentclass[aspectratio=169,10pt,english]{beamer}

\input{../../../_preambulo_beamer.tex}
\input{../../../_comandos_beamer.tex}

\selectlanguage{english}

\title{MA1001B - Statistical Modeling for Decision Making}
\subtitle{Lesson 45: Tukey HSD Multiple Comparisons}
\author{}
\institute{Tecnol\'ogico de Monterrey}
\date{}

\begin{document}

\begin{frame}[plain]
  \vspace{1.2cm}
  {\Large\bfseries MA1001B - Statistical Modeling for Decision Making\par}
  \vspace{0.7cm}
  {\large Lesson 45: Tukey HSD Multiple Comparisons\par}
  \vspace{0.7cm}
  {\color{TecRojo}\rule{\textwidth}{1.2pt}}
  \vspace{0.8cm}

  MA1001B Faculty\\[0.3cm]
  Term\\[0.3cm]
  Tecnol\'ogico de Monterrey
\end{frame}

\begin{frame}{The Pairwise Question After $F$}
  \begin{alertblock}{Guiding question}
    After a significant three-group $F$, which packing-method pairs differ,
    and how can that family of tests keep a 5 percent familywise error?
  \end{alertblock}

  \vspace{0.6em}
  By the end of this lesson, you can:
  \begin{itemize}
    \item run the three pairwise $t$ tests;
    \item read Tukey adjusted $p$-values and simultaneous CIs;
    \item name the pairs that differ at familywise $\alpha=0.05$;
    \item explain why unadjusted tests inflate familywise error.
  \end{itemize}

  \vfill
  \scriptsize All cycle times used today are synthetic. No real company data are used.
\end{frame}

\begin{frame}{Three Pairs, One Family}
  \small
  \begin{center}
    \begin{tabular}{lrr}
      \toprule
      \textbf{Pair} & \textbf{Mean difference} & \textbf{Unadjusted $p$} \\
      \midrule
      Standard $-$ Guided & 1.215833 & 0.287762 \\
      Standard $-$ Automated & 5.716390 & 0.000071 \\
      Guided $-$ Automated & 4.500557 & 0.000233 \\
      \bottomrule
    \end{tabular}
  \end{center}

  \vspace{0.5em}
  \begin{exampleblock}{Familywise error}
    Three tests at $\alpha=0.05$ are not one 5 percent decision. Familywise
    error is the chance of at least one false pair in the family.
  \end{exampleblock}
\end{frame}

\begin{frame}[fragile]{Step 1: Unadjusted Pairwise $t$}
  \begin{columns}[T]
    \begin{column}{0.40\textwidth}
      \small
      \begin{lstlisting}
t, p = ttest_ind(a, b,
                 equal_var=True)
reject = p < 0.05
      \end{lstlisting}

      \begin{block}{Verified output}
        Guided pair $p=0.287762$ retain\\
        Auto pairs $p=0.000071$, $0.000233$\\
        Three tests, no correction
      \end{block}
    \end{column}
    \begin{column}{0.60\textwidth}
      \centering
      \includegraphics[width=\textwidth]{../figures/tukey_hsd_01_pairwise_unadjusted.png}
      \vspace{0.2em}
      \scriptsize Unadjusted pairwise differences from Step 1
    \end{column}
  \end{columns}
\end{frame}

\begin{frame}{Tukey HSD Keeps the Family at 5 Percent}
  \small
  \begin{center}
    \begin{tabular}{lrrr}
      \toprule
      \textbf{Pair} & $\mathbf{p_{\text{adj}}}$ & \textbf{Tukey CI} & \textbf{Reject?} \\
      \midrule
      Auto vs Guided & 0.000789 & $(1.79,\ 7.21)$ & Yes \\
      Auto vs Standard & 0.000033 & $(3.00,\ 8.43)$ & Yes \\
      Guided vs Standard & 0.521330 & $(-1.50,\ 3.93)$ & No \\
      \bottomrule
    \end{tabular}
  \end{center}

  \vspace{0.4em}
  Automated differs from both other methods. Guided and Standard are not
  distinguished. The CI that contains 0 is the retain decision.
\end{frame}

\begin{frame}[fragile]{Step 2: Tukey Simultaneous Intervals}
  \begin{columns}[T]
    \begin{column}{0.40\textwidth}
      \small
      \begin{lstlisting}
tukey = pairwise_tukeyhsd(
  y, groups, alpha=0.05
)
      \end{lstlisting}

      \begin{block}{Verified output}
        Guided vs Standard retains\\
        $p_{\text{adj}}=0.521330$\\
        CI $(-1.498,\ 3.930)$
      \end{block}
    \end{column}
    \begin{column}{0.60\textwidth}
      \centering
      \includegraphics[width=\textwidth]{../figures/tukey_hsd_02_tukey_hsd.png}
      \vspace{0.2em}
      \scriptsize Tukey HSD intervals from Step 2
    \end{column}
  \end{columns}
\end{frame}

\begin{frame}{Step 3: Unadjusted Tests Inflate Familywise Error}
  \begin{columns}[T]
    \begin{column}{0.42\textwidth}
      \small
      5000 null experiments, all true means equal.

      \vspace{0.4em}
      Unadjusted pairwise FWER $=0.127800$.

      \vspace{0.4em}
      Tukey HSD FWER $=0.049000$ versus $\alpha=0.05$.

      \vspace{0.5em}
      \textbf{Limit:} pairwise $t$ tests without correction inflate
      familywise error. Lesson 46 checks the equal-variance assumption
      both procedures still need.
    \end{column}
    \begin{column}{0.58\textwidth}
      \centering
      \includegraphics[width=\textwidth]{../figures/tukey_hsd_03_familywise_error.png}
      \vspace{0.2em}
      \scriptsize Empirical familywise error from Step 3
    \end{column}
  \end{columns}
\end{frame}

\begin{frame}{Concept Check}
  \begin{enumerate}
    \item How many pairwise tests exist for $k=3$ methods?
    \item Why is familywise error not the same as one test's $\alpha=0.05$?
    \item Which pair is retained by Tukey HSD in this packing experiment?
    \item In 5000 null experiments, why is $0.127800$ a problem?
  \end{enumerate}

  \vspace{0.6em}
  \begin{block}{Connection to Lesson 46}
    The session can continue directly into ANOVA residual plots, Levene's
    test, and variance inequality.
  \end{block}

  \vfill
  \scriptsize Source: OpenStax, \textit{Introductory Business Statistics 2e},
  Sections 12.2--12.3, post-hoc pairwise practice. CC BY-NC-SA.
\end{frame}

\appendix



\end{document}
